% =====================================================================
%  2.2 Atividade com Git e GitHub
% =====================================================================
\subsection{Atividade com Git e GitHub}\label{sec:git}

\subsubsection{O que é Git}

O \textbf{Git} é um sistema de controle de versão \textbf{distribuído}: ele gerencia as mudanças
feitas nos arquivos de um projeto e mantém um histórico completo de modificações, registrando
\textbf{quem} alterou, \textbf{quando} alterou e \textbf{o que} foi alterado \citep{chacon2014}.
Uma boa analogia, usada em aula, é a do \textit{snapshot} de máquina virtual: cada versão registrada
é uma ``fotografia'' do estado do projeto naquele momento, o que permite voltar no tempo caso um
erro seja introduzido \citep{muniz2026}. Com isso o Git oferece rastreabilidade (é possível saber
qual linha mudou e por quê), evita a perda de código e permite que várias pessoas trabalhem no mesmo
projeto sem sobrescrever o trabalho umas das outras. O \textbf{GitHub}, por sua vez, é uma plataforma
\textit{on-line} que hospeda repositórios Git remotos e acrescenta recursos de colaboração, como
\textit{pull requests}, \textit{issues} e automações.

\subsubsection{Repositório local e repositório remoto}

O \textbf{repositório local} fica no próprio computador, dentro da pasta oculta \texttt{.git}
criada pelo \cod|git init| (ou pelo \cod|git clone|). Ele guarda todo o histórico do projeto, de
modo que é possível criar versões, consultar o histórico e criar \textit{branches} mesmo sem internet.
Como foi lembrado em aula, ``o Git mora dentro da pasta \texttt{.git}: se ela some, o Git também''.
O \textbf{repositório remoto} é uma cópia hospedada em um servidor -- no nosso caso, o GitHub --,
identificada normalmente pelo nome \texttt{origin}. Ele funciona como ponto central de sincronização
da equipe e como \textit{backup}: como cada desenvolvedor tem uma cópia completa do repositório,
mesmo que o servidor fique indisponível o código não se perde. A Figura~\ref{fig:fluxo-git} mostra
as áreas do Git e os comandos que movem as alterações entre elas.

\begin{figure}[H]
  \centering
  \caption{Áreas do Git e comandos que movem as alterações entre elas}\label{fig:fluxo-git}
  \begin{tikzpicture}
    \tikzset{area/.style={caixa, text width=2.25cm, minimum height=1.25cm}}
    \node[area] (wd) {\textbf{Diretório de trabalho}\\{\scriptsize arquivos editados}};
    \node[area, right=1.65cm of wd] (st) {\textbf{Área de preparação}\\{\scriptsize \textit{staging} (índice)}};
    \node[area, right=1.65cm of st] (lr) {\textbf{Repositório local}\\{\scriptsize pasta \texttt{.git}}};
    \node[caixalaranja, text width=2.25cm, minimum height=1.25cm, right=1.65cm of lr] (rr)
         {\textbf{Repositório remoto}\\{\scriptsize GitHub (\texttt{origin})}};
    \draw[seta] (wd.east) -- node[rotulo, above=1pt] {\texttt{git add}} (st.west);
    \draw[seta] (st.east) -- node[rotulo, above=1pt] {\texttt{git commit}} (lr.west);
    \draw[seta] ([yshift=4pt]lr.east) -- node[rotulo, above=1pt] {\texttt{git push}} ([yshift=4pt]rr.west);
    \draw[seta] ([yshift=-4pt]rr.west) -- node[rotulo, below=1pt] {\texttt{git fetch}} ([yshift=-4pt]lr.east);
    \draw[seta] (rr.south) -- ++(0,-0.75) -| node[rotulo, pos=0.25, below=1pt]
         {\texttt{git pull} / \texttt{git clone}: trazem as alterações do remoto para o diretório de trabalho} (wd.south);
    \begin{scope}[on background layer]
      \node[draw=corprimaria!45, dashed, rounded corners=4pt, fit=(wd)(st)(lr), inner sep=7pt,
            label={[font=\scriptsize\bfseries, text=corprimaria]above:Computador do desenvolvedor (repositório local)}] {};
    \end{scope}
  \end{tikzpicture}
  \fonte{Elaborado pelo autor (2026), com base em Chacon e Straub (2014).}
\end{figure}

\subsubsection{Commit, push, pull e branch}

\begin{itemize}
  \item \textbf{Commit}: registra no repositório local uma nova versão com as alterações que estavam
        na área de preparação. Cada \textit{commit} tem um identificador único (\textit{hash}), autor,
        data e uma mensagem que explica o que foi feito, por exemplo
        \cod|git commit -m "Adiciona playbook do app01"|.
  \item \textbf{Push}: envia os \textit{commits} do repositório local para o repositório remoto
        (\cod|git push|), tornando-os visíveis para a equipe.
  \item \textbf{Pull}: busca as alterações do repositório remoto e as integra à \textit{branch} local
        (\cod|git pull| equivale a um \cod|git fetch| seguido de um \cod|git merge|).
  \item \textbf{Branch}: é uma ramificação, ou seja, uma linha de desenvolvimento separada da principal
        (\texttt{main}). Permite desenvolver uma funcionalidade ou corrigir um erro de forma isolada,
        com \textit{commits} próprios, sem quebrar o código principal. Quando o trabalho termina, a
        \textit{branch} é integrada de volta com \cod|git merge|.
\end{itemize}

Na disciplina também foi apresentado o \textbf{Git Flow}, um modelo de organização de
\textit{branches} com cinco tipos: \texttt{main} (versões estáveis em produção, marcadas com
\textit{tags} como \texttt{v1.0}), \texttt{develop} (integração do desenvolvimento),
\texttt{feature} (novas funcionalidades, criadas a partir da \texttt{develop}), \texttt{release}
(preparação de uma nova versão, com testes finais) e \texttt{hotfix} (correções urgentes que nascem
da \texttt{main} e voltam para a \texttt{main} e para a \texttt{develop}). O Git Flow é indicado para
projetos grandes, com vários ambientes e lançamentos programados; já o \textit{feature branching},
mais simples, combina melhor com projetos menores e com integração contínua \citep{muniz2026}.

\subsubsection{Como o Git foi utilizado nas atividades}

Primeiro, o Git foi instalado e configurado com a identidade do autor, que fica gravada em cada
\textit{commit}. Em seguida, foi criado um repositório no GitHub e o projeto local foi inicializado
e ligado a ele, seguindo a sequência do Código~\ref{cod:git-sequencia}. Antes do primeiro
\textit{commit} foi criado o arquivo \texttt{.gitignore} (Código~\ref{cod:gitignore}), para que
arquivos que não devem ser publicados -- como a pasta \texttt{.vagrant/}, que guarda o estado
local das VMs, logs e arquivos \texttt{.env} com segredos -- ficassem fora do repositório.

\begin{terminal}[caption={Sequência de comandos usada para versionar o laboratório no GitHub},label={cod:git-sequencia}]
git --version                                   # confirma a instalação
git config --global user.name  "Seu Nome"       # identidade gravada nos commits
git config --global user.email "seu@email.com"
git config --global --list                      # confere a configuração

cd lab-devops                                   # ATENÇÃO: dentro da pasta do projeto
git init                                        # cria o repositório local (.git)
echo "# Lab DevOps" > README.md
git status                                      # mostra arquivos novos/alterados
git add .                                       # adiciona tudo à área de preparação
git commit -m "Primeiro commit do projeto"      # registra a primeira versão
git branch -M main                              # renomeia a branch principal para main
git remote add origin https://github.com/usuario/lab-devops.git
git remote -v                                   # confere o repositório remoto
git push -u origin main                         # envia os commits para o GitHub
\end{terminal}

% Conteúdo do arquivo codigos/.gitignore (inserido aqui porque o Overleaf pode ignorar arquivos ocultos)
\begin{listagem}[language=ini,caption={[Arquivo .gitignore do laboratório]{Arquivo .gitignore do laboratório -- \nolinkurl{.gitignore}}},label={cod:gitignore}]
# Arquivos que NÃO devem ir para o GitHub
.vagrant/
*.log
.env
.DS_Store
\end{listagem}

A partir daí, a cada etapa concluída (por exemplo, um novo \textit{playbook} ou um ajuste no
\texttt{Vagrantfile}) o ciclo \cod|git status| $\rightarrow$ \cod|git add| $\rightarrow$
\cod|git commit| $\rightarrow$ \cod|git push| foi repetido. Para alterações maiores foi usada uma
\textit{branch} separada, integrada à \texttt{main} somente depois de testada
(Código~\ref{cod:git-branch}). Este próprio relatório, junto com todos os arquivos do laboratório,
está versionado no repositório \url{\repositorio}.

\begin{terminal}[caption={Uso de branch para desenvolver uma alteração isolada},label={cod:git-branch}]
git checkout -b feature/playbook-app            # cria a branch e muda para ela
git add ansible/playbooks/app.yml
git commit -m "Adiciona playbook de configuração do app01"
git checkout main                               # volta para a branch principal
git pull                                        # garante que a main está atualizada
git merge feature/playbook-app                  # integra a alteração
git push                                        # publica no GitHub
git log --oneline --graph                       # histórico resumido
\end{terminal}

\subsubsection{Comandos utilizados}

\begin{table}[H]
  \centering
  \caption{Comandos do Git utilizados}\label{tab:comandos-git}
  \small
  \renewcommand{\arraystretch}{1.2}
  \begin{tabularx}{\linewidth}{>{\ttfamily}l >{\raggedright\arraybackslash}X}
    \toprule
    \normalfont\bfseries Comando & \textbf{Função} \\
    \midrule
    git init & Cria um repositório Git na pasta atual. \\
    git clone URL & Copia um repositório remoto (com todo o histórico) para o computador. \\
    git status & Mostra arquivos novos, alterados e preparados para \textit{commit}. \\
    git add . & Adiciona as alterações à área de preparação (\textit{staging}). \\
    git commit -m "msg" & Salva uma nova versão no repositório local. \\
    git remote add origin URL & Liga o repositório local ao remoto (GitHub). \\
    git pull & Busca e integra as alterações do remoto. \\
    git push & Envia os \textit{commits} locais para o remoto. \\
    git branch & Lista ou cria \textit{branches}. \\
    git checkout -b nome & Cria uma \textit{branch} e muda para ela. \\
    git merge nome & Integra a \textit{branch} indicada na \textit{branch} atual. \\
    git log -{}-oneline & Mostra o histórico de \textit{commits} de forma resumida. \\
    \bottomrule
  \end{tabularx}
  \fonte{Elaborado pelo autor (2026), com base em Muniz (2026).}
\end{table}

\subsubsection{Evidências da execução}

\printtela{github-repositorio}{Repositório do trabalho no GitHub}%
  {Página do repositório no GitHub mostrando os arquivos e pastas versionados
   (por exemplo, \texttt{relatorio/} e \texttt{README.md}) e o último \textit{commit}.}

\printtela{git-status}{Saída do comando \texttt{git status}}%
  {Terminal na pasta do projeto mostrando o \texttt{git status} (arquivos modificados/não rastreados
   ou ``nothing to commit, working tree clean'').}

\printtela{git-log}{Histórico de commits (\texttt{git log -{}-oneline})}%
  {Terminal com \texttt{git log -{}-oneline -{}-graph} mostrando os \textit{commits} realizados.}

\printtela{git-push}{Envio das alterações para o GitHub com \texttt{git push}}%
  {Terminal com a saída do \texttt{git push} (ou a aba \textit{Commits} do GitHub com o histórico).}

\printtela{github-arquivo}{Arquivo versionado visualizado no GitHub}%
  {GitHub exibindo um arquivo versionado, por exemplo o \texttt{Vagrantfile} ou o
   \texttt{playbooks/db.yml}, com o histórico de alterações.}

\subsubsection{Por que o versionamento é importante em DevOps}

No DevOps, o repositório Git é a ``fonte única da verdade'': tudo o que define o sistema -- o código
da aplicação, os \texttt{Vagrantfile}, os \textit{playbooks} do Ansible, o \texttt{Dockerfile} e o
\texttt{Jenkinsfile} -- fica versionado. Isso traz rastreabilidade (cada mudança tem autor, data e
motivo), segurança para experimentar (é possível voltar a uma versão que funcionava), colaboração
organizada por meio de \textit{branches} e revisões, e é a base da automação: os \textit{pipelines}
de CI/CD são disparados justamente por um \textit{commit} ou \textit{push} no repositório. Sem
versionamento, não existe Infraestrutura como Código nem integração contínua.
